% ============================================================
\chapter{Scott Domains and Scott Topology}
\label{ch:scott}
% ============================================================

In 1970, Dana Scott was searching for a mathematical model for the untyped lambda calculus---a formal system where a function can take itself as an argument. Classical topological spaces would not do: what was needed was a space isomorphic to its own function space. Scott discovered that certain complete partial orders, equipped with a carefully chosen topology---the \emph{Scott topology}---possess exactly this property. \emph{Scott domains} thus became the foundation of denotational semantics for programming languages, revealing a deep connection between topology, logic, and computer science.

\begin{intuition}
The Scott topology is the natural topology on a directed-complete partial order (dcpo) that
captures the notion of ``upward convergence.'' It is the fundamental topology in
denotational semantics, where programs are modeled as elements of domains and the
denotation of a recursive program is the supremum of a chain of approximations.
\end{intuition}

% ------------------------------------------------------------
\section{Directed-Complete Partial Orders}
% ------------------------------------------------------------

\begin{definition}[Directed set]
Let $(P,\leq)$ be a partially ordered set. A subset $D \subseteq P$ is \emph{directed}
if $D \neq \emptyset$ and for all $d_1, d_2 \in D$, there exists $d \in D$ with
$d_1 \leq d$ and $d_2 \leq d$.
\end{definition}

\begin{definition}[Dcpo]
A \emph{dcpo} (directed-complete partial order) is a poset $(D,\leq)$ in which every
directed subset has a supremum. We write $\bigsqcup S$ or $\bigvee S$ for the supremum of
a directed set $S$.
\end{definition}

\begin{example}[Examples of dcpos]
\begin{enumerate}
  \item Every complete lattice is a dcpo.
  \item $(\mathcal{P}(X), \subseteq)$ is a dcpo (and a complete lattice).
  \item $([0,1], \leq)$ is a dcpo.
  \item The set $\Sigma^{\leq\omega}$ of finite and infinite words over an alphabet
    $\Sigma$, ordered by the prefix relation, is a dcpo.
  \item The set $X \rightharpoonup Y$ of partial functions from $X$ to $Y$, ordered by
    graph extension, is a dcpo.
\end{enumerate}
\end{example}

\begin{definition}[Pointed dcpo]
A dcpo is \emph{pointed} if it has a least element, denoted $\bot$.
\end{definition}

% ------------------------------------------------------------
\section{The Scott Topology}
% ------------------------------------------------------------

\begin{definition}[Scott topology]
Let $(D,\leq)$ be a dcpo. A subset $U \subseteq D$ is \emph{Scott-open} if:
\begin{enumerate}
  \item $U$ is upward-closed: $x \in U$ and $x \leq y \Rightarrow y \in U$.
  \item $U$ is \emph{inaccessible by directed suprema}: if $S$ is directed and
    $\bigsqcup S \in U$, then $S \cap U \neq \emptyset$.
\end{enumerate}
The collection of Scott-open sets forms the \emph{Scott topology} $\sigma(D)$.
\end{definition}

\begin{proposition}[The Scott topology is indeed a topology]
The Scott-open sets satisfy:
\begin{enumerate}
  \item $\emptyset$ and $D$ are Scott-open.
  \item Any union of Scott-open sets is Scott-open.
  \item Any finite intersection of Scott-open sets is Scott-open.
\end{enumerate}
\end{proposition}

\begin{proof}
Only (3) requires verification. Let $U_1, U_2$ be Scott-open and $S$ directed with
$\bigsqcup S \in U_1 \cap U_2$. There exist $s_1 \in S \cap U_1$ and
$s_2 \in S \cap U_2$. By directedness, there exists $s \in S$ with $s_1 \leq s$ and
$s_2 \leq s$. Since $U_1$ and $U_2$ are upward-closed, $s \in U_1 \cap U_2$, so
$S \cap (U_1 \cap U_2) \neq \emptyset$.
\end{proof}

\begin{proposition}[Specialization order]
The specialization order of the Scott topology coincides with the original order:
$x \leq_\sigma y \Leftrightarrow x \leq y$.
\end{proposition}

% ------------------------------------------------------------
\section{Scott-Continuous Functions}
% ------------------------------------------------------------

\begin{definition}[Scott-continuous function]
Let $D$ and $E$ be dcpos. A function $f\colon D \to E$ is \emph{Scott-continuous} if it is
continuous for the Scott topologies, equivalently:
\begin{enumerate}
  \item $f$ is monotone: $x \leq y \Rightarrow f(x) \leq f(y)$.
  \item $f$ preserves directed suprema: for every directed set $S$,
    $f(\bigsqcup S) = \bigsqcup f(S)$.
\end{enumerate}
\end{definition}

\begin{theorem}[Characterization of Scott-continuity]
For a function $f\colon D \to E$ between dcpos, the following are equivalent:
\begin{enumerate}
  \item $f$ is continuous for the Scott topologies.
  \item $f$ is monotone and preserves directed suprema.
\end{enumerate}
\end{theorem}

\begin{proof}
$(1) \Rightarrow (2)$: Continuity implies monotonicity (since the specialization order is
the original order). Let $S$ be directed in $D$. Since $f$ is monotone, $f(\bigsqcup S)$
is an upper bound of $f(S)$. Let $e$ be any upper bound of $f(S)$ and suppose
$f(\bigsqcup S) \not\leq e$. Then there is a Scott-open $V \ni f(\bigsqcup S)$ with
$e \notin V$. The preimage $f^{-1}(V)$ is Scott-open and contains $\bigsqcup S$, so some
$s \in S$ is in $f^{-1}(V)$, giving $f(s) \in V$. Since $f(s) \leq e$ and $V$ is
upward-closed, $e \in V$, contradiction. So $f(\bigsqcup S) = \bigsqcup f(S)$.

$(2) \Rightarrow (1)$: Let $V$ be Scott-open in $E$. Then $f^{-1}(V)$ is upward-closed
since $f$ is monotone. If $S$ is directed with $\bigsqcup S \in f^{-1}(V)$, then
$f(\bigsqcup S) = \bigsqcup f(S) \in V$. Since $V$ is Scott-open and $f(S)$ is directed,
some $s \in S$ has $f(s) \in V$, so $s \in f^{-1}(V)$.
\end{proof}

% ------------------------------------------------------------
\section{The Category $\mathbf{Dcpo}$}
% ------------------------------------------------------------

\begin{theorem}[Properties of $\mathbf{Dcpo}$]
\begin{enumerate}
  \item The category $\mathbf{Dcpo}$ of dcpos and Scott-continuous functions is cartesian
    closed.
  \item The product of dcpos is the cartesian product with the componentwise order.
  \item The exponential $[D \to E]$ is the set of Scott-continuous functions from $D$ to
    $E$, ordered pointwise.
  \item $[D \to E]$ is a dcpo, and the Scott topology on $[D \to E]$ is the Isbell
    topology (which coincides with the compact-open topology for compact Scott-open sets).
\end{enumerate}
\end{theorem}

\begin{proof}[Proof sketch]
For (3): let $(f_i)_{i\in I}$ be directed in $[D \to E]$. For every $d \in D$,
$(f_i(d))_{i\in I}$ is directed in $E$ (by pointwise monotonicity and directedness of
$(f_i)$). Define $(\bigsqcup f_i)(d) = \bigsqcup_i f_i(d)$. One must verify that
$\bigsqcup f_i$ is Scott-continuous, which uses the fact that directed suprema commute in
a dcpo.
\end{proof}

% ------------------------------------------------------------
\section{Continuous Domains}
% ------------------------------------------------------------

\begin{definition}[Way-below relation]
In a dcpo $D$, we say $x$ is \emph{way below} $y$ (written $x \ll y$) if for every
directed set $S$ with $y \leq \bigsqcup S$, there exists $s \in S$ with $x \leq s$.
\end{definition}

\begin{definition}[Basis of a dcpo]
A subset $B \subseteq D$ is a \emph{basis} for the dcpo $D$ if for every $x \in D$, the
set $\{b \in B : b \ll x\}$ is directed with supremum $x$.
\end{definition}

\begin{definition}[Continuous domain]
A dcpo $D$ is a \emph{continuous domain} if for every $x \in D$, the set
$\mathop{\Downarrow}\! x = \{y \in D : y \ll x\}$ is directed and
$x = \bigsqcup(\mathop{\Downarrow}\! x)$.
\end{definition}

\begin{theorem}[Scott topology of a continuous domain]
If $D$ is a continuous domain, then:
\begin{enumerate}
  \item The sets $\mathop{\uparrow\!\!\uparrow} b = \{x \in D : b \ll x\}$ form a base
    for the Scott topology.
  \item The Scott topology is sober.
  \item The lattice $\sigma(D)$ is a continuous lattice.
  \item $(D,\sigma(D))$ is locally compact.
\end{enumerate}
\end{theorem}

\begin{proof}
(1) Let $U$ be Scott-open and $x \in U$. Since $x = \bigsqcup(\mathop{\Downarrow}\! x)$
and $U$ is inaccessible, there exists $b \ll x$ with $b \in U$. Then
$x \in \mathop{\uparrow\!\!\uparrow} b \subseteq \uparrow\! b \subseteq U$ since $U$ is
upward-closed.

(2) Let $F$ be an irreducible closed set. The set
$B_F = \{b : \exists x \in F,\, b \ll x\}$ is directed. Indeed, if $b_1 \ll x_1 \in F$
and $b_2 \ll x_2 \in F$, then $\mathop{\uparrow\!\!\uparrow} b_1$ and
$\mathop{\uparrow\!\!\uparrow} b_2$ are opens meeting $F$, so by irreducibility,
$F \cap \mathop{\uparrow\!\!\uparrow} b_1 \cap \mathop{\uparrow\!\!\uparrow} b_2 \neq
\emptyset$, giving some $y \in F$ with $b_1 \ll y$ and $b_2 \ll y$. The supremum
$\bigsqcup B_F$ lies in $F$ (since $F$ is Scott-closed) and is the generic point.
\end{proof}

% ------------------------------------------------------------
\section{Algebraic Domains}
% ------------------------------------------------------------

\begin{definition}[Compact element]
An element $k \in D$ is \emph{compact} (or \emph{finite}) if $k \ll k$. The set of
compact elements is denoted $K(D)$.
\end{definition}

\begin{definition}[Algebraic domain]
A dcpo $D$ is \emph{algebraic} if for every $x \in D$, the set
$\{k \in K(D) : k \leq x\}$ is directed and $x = \bigsqcup\{k \in K(D) : k \leq x\}$.
\end{definition}

\begin{proposition}
Every algebraic domain is a continuous domain. The converse is false.
\end{proposition}

\begin{example}
$(\mathcal{P}(\N), \subseteq)$ is an algebraic domain: the compact elements are the finite
subsets. The interval domain $\mathbf{I}\R = \{[a,b] : a \leq b, a,b \in \R\}$ ordered by
reverse inclusion is a continuous domain that is not algebraic.
\end{example}

% ------------------------------------------------------------
\section{Representation Theorem}
% ------------------------------------------------------------

\begin{theorem}[Algebraic domains and ideals]
If $D$ is an algebraic domain, then $D$ is isomorphic (as a dcpo) to the set
$\mathrm{Idl}(K(D))$ of ideals of $(K(D), \leq)$ (directed downward-closed subsets of
$K(D)$), ordered by inclusion.
\end{theorem}

\begin{proof}
The isomorphism sends $x \in D$ to $\downarrow\! x \cap K(D)$. This is an ideal: it is
downward-closed in $K(D)$ by definition, and directed because $D$ is algebraic. The
inverse sends an ideal $I$ to $\bigsqcup I$. The key point is that directed unions of
ideals correspond to directed suprema in $D$.
\end{proof}

% ------------------------------------------------------------
\section{Exercises}
% ------------------------------------------------------------

\begin{exercise}
Show that the Scott topology on $(\R, \leq)$ consists of $\emptyset$, $\R$, and the
intervals $(a, +\infty)$ for $a \in \R$.
\end{exercise}

\begin{exercise}
Show that the Scott topology on $(\mathcal{P}(\N), \subseteq)$ is not metrizable.
\end{exercise}

\begin{exercise}
Show that the product of two dcpos (with the product order) is a dcpo and that the Scott
topology of the product is finer than the product of the Scott topologies. Give an example
where the inclusion is strict.
\end{exercise}

\begin{exercise}
Let $D$ be a pointed dcpo. Show that the function $\mathrm{fix}\colon [D \to D] \to D$
sending $f$ to its least fixed point is Scott-continuous.
\end{exercise}

\begin{exercise}
Let $D$ be a continuous domain. Show that $x \ll y$ if and only if
$y \in \mathrm{Int}(\uparrow\! x)$ (Scott interior).
\end{exercise}

\begin{exercise}
Show that a dcpo $D$ is continuous if and only if $\sigma(D)$ is a completely distributive
lattice.
\end{exercise}

\begin{exercise}[Step functions]
Let $D$ be an algebraic domain and $E$ a dcpo. Show that every Scott-continuous function
$f\colon D \to E$ is the directed supremum of ``step functions'' (functions defined on
compact elements of $D$ and extended by directed suprema).
\end{exercise}
